\section{Theoretical introduction}
\nocite{TietzeSchenk, TL06x}
An operational amplifier (op-amp) is a differential amplifier with a very high voltage gain, a high input impedance and a low output impedance.
In the circuits below it always works with negative feedback, so its output settles at the value for which the voltage between the inputs is close to zero.
Together with the negligible input currents this gives the two rules of the ideal op-amp used throughout: $U_+ = U_-$ and $I_+ = I_- = 0$.

The op-amp used in this laboratory is the TL061, a low-power JFET-input op-amp.
Its parameters relevant here are listed in \autoref{tab:tl061}.
\begin{table}[ht]
  \centering
  \caption{TL061C parameters at $V_{CC\pm} = \pm\qty{15}{V}$ and \qty{25}{\celsius}~\cite{TL06x}}
  \label{tab:tl061}
  \begin{tabular}{lrr}
    \hline
    Parameter & Typical & Worst case \\
    \hline
    Input offset voltage $V_{IO}$ & \qty{3}{mV} & \qty{15}{mV} \\
    Input bias current $I_{IB}$ & \qty{30}{pA} & \qty{400}{pA} \\
    Unity-gain bandwidth $f_T$ & \qty{1}{MHz} & \\
    Slew rate $SR$ & \qty{3.5}{V/\us} & \qty{1.5}{V/\us} \\
    Common-mode rejection ratio & \qty{86}{dB} & \qty{70}{dB} \\
    \hline
  \end{tabular}
\end{table}

\subsection{Inverting amplifier}\label{sec:th_inverting}
The input signal is applied through $R_1$ to the inverting input, and $R_4$ closes the feedback loop.
The inverting input is a virtual ground, so the current $U_{in}/R_1$ flows through $R_4$ and
\begin{equation}\label{eqn:inverting}
  G = \frac{U_{out}}{U_{in}} = -\frac{R_4}{R_1}.
\end{equation}
The bandwidth is limited by the gain-bandwidth product of the op-amp: the upper cut-off frequency is $f_H \approx f_T / (1 + R_4/R_1)$.
A capacitor $C_3$ in parallel with $R_4$ adds a pole at $f_c = 1/(2\pi R_4 C_3)$.

\subsection{Non-inverting amplifier}\label{sec:th_noninverting}
The signal is applied to the non-inverting input, and $R_1$ connects the inverting input to ground.
\begin{equation}\label{eqn:noninverting}
  G = 1 + \frac{R_4}{R_1}, \qquad f_H \approx \frac{f_T}{G}.
\end{equation}
With $C_3$ in parallel with $R_4$, the gain falls from $G$ above $1/(2\pi R_4 C_3)$, but only down to unity.

\subsection{Integrator}\label{sec:th_integrator}
With a capacitor $C_3$ in the feedback path, the current $U_{in}/R_1$ charges $C_3$:
\begin{equation}\label{eqn:integrator}
  U_{out}(t) = -\frac{1}{R_1 C_3}\int U_{in}\,dt, \qquad H(j\omega) = -\frac{1}{j\omega R_1 C_3}.
\end{equation}
A square wave of amplitude $U_{in}$ and frequency $f$ is integrated into a triangle wave of amplitude $U_{out} = U_{in}/(4 f R_1 C_3)$.
The resistor $R_4$ in parallel with $C_3$ limits the DC gain to $-R_4/R_1$, so the circuit integrates only for $f \gg 1/(2\pi R_4 C_3)$.

\subsection{Differentiator}\label{sec:th_differentiator}
With a capacitor $C_1$ at the input and $R_4$ in the feedback path:
\begin{equation}\label{eqn:differentiator}
  U_{out}(t) = -R_4 C_1 \frac{dU_{in}}{dt}, \qquad H(j\omega) = -j\omega R_4 C_1.
\end{equation}
A triangle wave of amplitude $U_{in}$ and frequency $f$ becomes a square wave of amplitude $U_{out} = 4 f R_4 C_1 U_{in}$.
The resistor $R_1$ in series with $C_1$ limits the gain at high frequencies and prevents ringing; the circuit differentiates up to about $1/(2\pi R_1 C_1)$.

\subsection{Non-ideal parameters}\label{sec:th_nonideal}
\begin{description}
  \item[Offset voltage] The DC voltage between the inputs needed to bring the output to zero. It appears at the output multiplied by the noise gain $1 + R_4/R_1$.
  \item[CMRR] The ratio of the differential to the common-mode gain, $\mathrm{CMRR} = 20\log_{10}(A_d/A_{cm})$.
  \item[Slew rate] The maximum rate of change of the output voltage. A sine of peak value $U_p$ is undistorted only up to $f = SR/(2\pi U_p)$.
\end{description}

\subsection{Transfer function}\label{sec:th_transfer}
The transfer function $H(j\omega) = U_{out}(j\omega)/U_{in}(j\omega)$ is measured with a sine input: at each frequency the magnitude is the ratio of the amplitudes, expressed in decibels as $K = 20\log_{10}(|U_{out}|/|U_{in}|)$, and the phase is $\varphi = 360^\circ \cdot f \cdot \Delta t$, where $\Delta t$ is the time shift between the input and the output.
